\daytitle{22}{洛伦兹力与带电粒子在磁场中的运动}

\begin{learningbox}{本节导读}
上一节的安培力研究的是磁场对一段电流的作用。电流来自大量电荷的定向移动，所以本节进一步研究磁场对单个运动电荷的作用：洛伦兹力。读题时先定 $v$、$B$ 和电荷正负，判断力的方向；再利用“磁场力不做功”判断速率；最后把轨迹圆弧转化为圆心、半径和圆心角的几何问题。
\end{learningbox}

\textbf{本节读题路线：} 先画出 $v$ 与 $B$，并按电荷正负判断洛伦兹力方向，确定粒子的偏转方向；再用“不做功”判断速率是否变化；最后把弯曲轨迹看成圆周运动的一段，用圆心、半径和圆心角解决几何问题。半径公式和周期公式都可以从受力图与圆周运动方程中得到。

\section*{一、学习目标}
\begin{enumerate}
  \item 能说明洛伦兹力是磁场对运动电荷的作用力。
  \item 能根据速度与磁场夹角使用洛伦兹力公式 $F=|q|vB\sin\theta$，并判断力的方向。
  \item 能解释磁场洛伦兹力为什么不做功，只改变速度方向、不改变速率。
  \item 能由洛伦兹力提供向心力推出匀强磁场中的半径 $r=\dfrac{mv}{|q|B}$ 和周期 $T=\dfrac{2\pi m}{|q|B}$。
  \item 会分析粒子进入有界磁场时的圆心、半径、偏转角和运动时间。
  \item 能说明回旋加速器中“电场加速、磁场转弯”的基本分工。
\end{enumerate}
\section*{二、初高中衔接}

Day 17 中，带电粒子在匀强电场中受到电场力，可能做类平抛运动。磁场中的情况很不同：
\begin{itemize}
  \item 电场力对正电荷沿电场方向、对负电荷反向，可能改变速率；
  \item 磁场力（高中常称洛伦兹力）始终垂直于速度方向，只改变速度方向；
  \item 电场力可以改变粒子动能，只有磁场力单独作用时，粒子速率不变。
\end{itemize}

因此，磁场题首先要判断“速率是否改变”，再研究“受力方向怎样使速度转向、怎样由轨迹确定圆心和半径”。

\textbf{科学发展视角：} 洛伦兹力是粒子加速器、质谱仪、等离子体约束和空间辐射探测的基础。现代科学常常不是直接“看见”微观粒子，而是通过电场和磁场控制它们的轨迹，再由轨迹反推粒子的速度、质量和电荷性质。

\section*{三、核心概念}

\begin{conceptbox}{洛伦兹力}

运动电荷在磁场中受到的力叫洛伦兹力。若电荷量为 $q$，速度为 $v$，磁感应强度为 $B$，速度方向与磁场方向夹角为 $\theta$，则
\[
\boxed{F = |q|vB\sin\theta}
\]

\begin{itemize}
  \item 当 $v\perp B$ 时，$F=|q|vB$。
  \item 当 $v$ 与 $B$ 平行或反平行时，$F=0$。
  \item 公式中的 $|q|$ 用电荷量大小；方向要另行判断。
\end{itemize}

\begin{center}
% TikZ 图：洛伦兹力方向
% Day 22：洛伦兹力
\begin{tikzpicture}[scale=0.92,>=Latex]
  \foreach \x in {-1.8,-1.0,-0.2,0.6,1.4} {
    \foreach \y in {-0.8,0,0.8} {
      \draw[blue,thick] (\x-0.08,\y-0.08) -- (\x+0.08,\y+0.08);
      \draw[blue,thick] (\x-0.08,\y+0.08) -- (\x+0.08,\y-0.08);
    }
  }
  \node[blue] at (2.15,0.85) {$B$ 向里};
  \fill[red] (-1.55,0) circle (2pt);
  \node[red,below] at (-1.55,-0.12) {$+q$};
  \draw[->,red,very thick] (-1.55,0) -- (-0.1,0) node[below] {$v$};
  \draw[->,purple,very thick] (-1.0,0.05) -- (-1.0,1.15) node[above] {$F$};
  \fill[red] (0.75,0) circle (2pt);
  \node[red,below] at (0.75,-0.12) {$-q$};
  \draw[->,red,very thick] (0.75,0) -- (1.85,0) node[below] {$v$};
  \draw[->,purple,very thick] (1.25,-0.05) -- (1.25,-1.15) node[below] {$F$};
\end{tikzpicture}


\figcaption{正电荷受力方向用左手定则判断；负电荷方向相反}\label{fig:lorentz-force-direction}
\end{center}

\textbf{方向判断：} 对正电荷，四指指向速度方向，磁感线穿入手心，拇指指向洛伦兹力方向。对负电荷，方向与正电荷相反。

\end{conceptbox}

\begin{conceptbox}{安培力的微观来源}

上一节的安培力和本节的洛伦兹力不是两个互不相关的力。通电导线中有大量载流子定向移动，磁场先作用在每个运动电荷上；载流子再通过与晶格的相互作用把力传给整段导线，宏观上就表现为安培力。为使速度方向与约定电流方向一致，下面采用\textbf{正载流子等效模型}推导力的大小。

\begin{center}
% TikZ 图：安培力的微观来源
% Day 22：洛伦兹力
\begin{tikzpicture}[scale=0.9,>=Latex,line cap=round]
  \path[use as bounding box] (-0.35,-2.25) rectangle (7.85,2.30);
  \tikzset{
    wire/.style={thick,fill=gray!12},
    carrier/.style={circle,fill=red!12,draw=red!80!black,inner sep=1.2pt},
    vec/.style={->,very thick}
  }

  \draw[wire] (0,0) rectangle (7.0,0.90);
  \draw[<->,thick] (0,1.05) -- (7.0,1.05);
  \node[above] at (3.50,1.05) {$L$};
  \node[font=\small] at (3.50,-0.28) {正载流子等效模型};

  \foreach \x in {0.65,1.45,...,6.65} {
    \foreach \y in {-0.75,-1.35} {
      \draw[blue!70!black,thick] (\x-0.09,\y-0.09) -- (\x+0.09,\y+0.09);
      \draw[blue!70!black,thick] (\x-0.09,\y+0.09) -- (\x+0.09,\y-0.09);
    }
  }
  \node[blue!70!black,font=\small] at (7.35,-1.05) {$B$ 向里};

  \foreach \x/\y in {0.75/0.28,1.75/0.62,2.75/0.28,3.75/0.62,4.75/0.28,5.75/0.62} {
    \node[carrier,font=\scriptsize] at (\x,\y) {$+$};
    \draw[vec,red] (\x+0.13,\y) -- ++(0.44,0);
  }
  \node[red,font=\small] at (3.45,0.36) {$v_d$};

  \draw[vec,red] (0.45,1.55) -- (2.10,1.55);
  \node[red,above] at (1.28,1.55) {$I=n|q|Sv_d$};

  \draw[vec,orange!85!black] (6.35,1.18) -- (6.35,1.82);
  \node[orange!85!black,right,font=\small] at (6.35,1.52)
    {$F_{\text{安}}=BIL$};

  \node[font=\small,align=center] at (3.50,-1.95)
    {载流子所受磁场力传给晶格，宏观上表现为导线所受安培力};
\end{tikzpicture}


\figcaption{正载流子等效模型：微观洛伦兹力经导体传递，宏观上表现为安培力}\label{fig:ampere-lorentz-micro}
\end{center}

以导线与匀强磁场垂直为例。设单位体积内载流子数为 $n$，每个载流子电荷量大小为 $|q|$，导线横截面积为 $S$，载流子定向移动速率为 $v_d$。长度为 $L$ 的导线内载流子数约为
\[
N=nSL.
\]
每个等效正载流子受到的洛伦兹力大小为 $|q|v_dB$。这些力传给晶格后，整段导线受到的安培力大小为
\[
F=N|q|v_dB=nSL|q|v_dB.
\]
又因为电流的微观表达式为
\[
I=n|q|Sv_d,
\]
于是
\[
F=BIL.
\]
若电流方向与磁场方向成角度 $\theta$，可以把载流子速度分解成平行磁场和垂直磁场两部分。平行磁场的分量不受磁场力，只有垂直分量 $v_d\sin\theta$ 产生作用，因此得到一般式
\[
F=BIL\sin\theta.
\]

\textbf{认知提醒：} 高中计算安培力时直接用约定电流方向和左手定则。金属中真实定向移动的是自由电子，电子漂移方向与约定电流方向相反；电子所受磁场力通过与晶格的相互作用传给导线，所得宏观安培力方向仍可直接用约定电流方向判断。

\end{conceptbox}

\begin{conceptbox}{洛伦兹力不做功}

本节所说洛伦兹力，指磁场对运动电荷的力。它的方向始终垂直于速度方向，因此瞬时功率
\[
P = \vect{F}\cdot \vect{v}=0
\]
所以磁场洛伦兹力不做功，不改变粒子动能和速率，只改变速度方向。

\begin{center}
% TikZ 图：洛伦兹力不做功
% Day 22：洛伦兹力
\begin{tikzpicture}[scale=0.95,>=Latex]
  \draw[thick,dashed] (0,0) circle (1.35);
  \fill[red] (0.95,0.95) circle (2pt);
  \draw[->,red,very thick] (0.95,0.95) -- (1.75,0.15) node[right] {$v$};
  \draw[->,purple,very thick] (0.95,0.95) -- (0.35,0.35) node[left] {$F$};
  \draw[gray] (0,0) -- (0.95,0.95);
  \node at (0,-1.85) {$F\perp v,\quad P=\vect{F}\cdot\vect{v}=0$};
\end{tikzpicture}


\figcaption{磁场洛伦兹力始终垂直速度：只改变方向，不改变速率}\label{fig:lorentz-no-work}
\end{center}

\textbf{重要区别：} 判断动能是否改变时先看力与速度是否有同向分量。电场力可能与速度同向或反向，可以改变带电粒子的动能；磁场洛伦兹力始终垂直速度，单独作用时只改变速度方向，不改变速率和动能。

\end{conceptbox}

\begin{conceptbox}{匀强磁场中的圆周运动}

若带电粒子以速度 $v$ 垂直进入匀强磁场，磁场洛伦兹力始终垂直速度，提供向心力：
\[
|q|vB=\frac{mv^2}{r}
\]
这一步要读成“洛伦兹力充当圆周运动所需的向心力”，不要理解成粒子同时又多受到一个叫向心力的新力。实际作用在粒子上的磁场力是 $|q|vB$；$\dfrac{mv^2}{r}$ 是匀速圆周运动所需合力大小的表达式。两者相等，表示磁场力的大小刚好满足转弯需要。

由向心力方程先得到轨道半径
\[
\boxed{r=\frac{mv}{|q|B}}.
\]
粒子运动一周的路程为 $2\pi r$，速率保持为 $v$，所以
\[
T=\frac{2\pi r}{v}
=\frac{2\pi}{v}\cdot\frac{mv}{|q|B}
=\boxed{\frac{2\pi m}{|q|B}}.
\]
推导中 $v$ 被约去，说明在非相对论条件下，同一种粒子垂直进入同一匀强磁场时，回旋周期与速率无关；速率增大时，轨道半径按相同比例增大。

\begin{center}
% TikZ 图：磁场中的圆周运动
% Day 22：洛伦兹力
\begin{tikzpicture}[scale=0.9,>=Latex]
  \foreach \x in {-2.0,-1.2,-0.4,0.4,1.2,2.0} {
    \foreach \y in {-1.3,-0.5,0.3,1.1} {
      \draw[blue,thick] (\x-0.07,\y-0.07) -- (\x+0.07,\y+0.07);
      \draw[blue,thick] (\x-0.07,\y+0.07) -- (\x+0.07,\y-0.07);
    }
  }
  \node[blue] at (2.65,1.1) {$B$ 向里};
  \draw[thick,dashed] (0,0) circle (1.15);
  \coordinate (P) at (0,-1.15);
  \draw[->,red,very thick] (-1.65,-1.15) -- (P) node[midway,below] {$v$};
  \draw[red,thick,->] (P) arc (-90:170:1.15);
  \fill[red] (P) circle (2pt);
  \fill (0,0) circle (1.5pt);
  \draw[gray] (0,0) -- (P) node[midway,right] {$r$};
  \draw[->,purple,thick] (P) -- (0,-0.35) node[midway,right] {$F$};
\end{tikzpicture}


\figcaption{带电粒子垂直进入匀强磁场后做匀速圆周运动}\label{fig:magnetic-circular-motion}
\end{center}

\begin{itemize}
  \item 半径 $r$ 与速度 $v$ 成正比。
  \item 在非相对论、匀强磁场且速度垂直磁场的圆周运动模型中，周期 $T$ 与速度 $v$ 无关，只由 $m$、$|q|$、$B$ 决定。
  \item 轨迹方向由电荷正负、速度方向和磁场方向共同决定。
\end{itemize}

\end{conceptbox}

\begin{samepage}
\section*{四、公式与适用条件}

\textbf{模型使用前检查：}
\nopagebreak[4]
\begin{enumerate}
  \item 先判断粒子电荷正负，负电荷的洛伦兹力方向与正电荷相反。
  \item $F=|q|vB$ 默认 $v\perp B$；若不垂直，应使用 $F=|q|vB\sin\theta$。
  \item 匀强磁场圆周运动中，先找圆心、半径和圆心角，再写时间关系。
  \item 电场和磁场同时存在时，先分别判断电场力和磁场力，再做合力分析。
\end{enumerate}
\end{samepage}

\begin{center}
\begingroup
\renewcommand{\arraystretch}{1.35}
\begin{tabular}{lll}
\toprule
公式 & 物理意义 & 适用条件 \\
\midrule
$F=|q|vB\sin\theta$ & 洛伦兹力大小 & 运动电荷在磁场中 \\
$F=|q|vB$ & 最大洛伦兹力 & $v\perp B$ \\
$F=BIL\sin\theta$ & 安培力宏观表达 & 直导线在匀强磁场中 \\
$r=\dfrac{mv}{|q|B}$ & 圆周运动半径 & $v\perp B$，匀强磁场 \\
$T=\dfrac{2\pi m}{|q|B}$ & 圆周运动周期 & $v\perp B$，匀强磁场 \\
$t=\dfrac{\alpha}{2\pi}T$ & 磁场中运动时间 & 圆心角为 $\alpha$（弧度制） \\
\bottomrule
\end{tabular}
\endgroup
\end{center}

\begin{learningbox}{洛伦兹力圆弧题的四步分析清单}
带电粒子进入磁场后，先完成方向和圆心判断，再代入半径公式。按下面四步走，几何关系会清楚很多：
\begin{enumerate}
  \item \textbf{定偏转方向：} 先看电荷正负，再由 $v$ 和 $B$ 判断洛伦兹力方向，确定轨迹偏向哪一侧。
  \item \textbf{找圆心：} 洛伦兹力始终指向圆心，所以圆心一定在速度的垂线上；入射点、出射点处都可以作速度垂线。
  \item \textbf{定半径和圆心角：} 半径来自 $r=\dfrac{mv}{|q|B}$ 或题图几何；圆心角等于速度方向转过的角度。
  \item \textbf{算时间或位置：} 时间用 $t=\dfrac{\alpha}{2\pi}T$，位置用弦长、半径和圆心角的几何关系。
\end{enumerate}
如果题目中同时有电场，先分清“电场改变速率、磁场改变方向”；只有磁场单独作用时，速率不变，圆周半径才保持不变。
\end{learningbox}

\section*{五、典型模型}

\textbf{磁场粒子题的三层拆解：} 遇到带电粒子进磁场，建议把半径公式放在“方向判断”和“圆周模型”之后使用。完整思路是：
\[
\text{方向判断}\rightarrow\text{圆周模型}\rightarrow\text{边界几何}.
\]
第一步确定粒子的偏转方向；第二步确认洛伦兹力提供向心力；第三步再看题目给的是半径、直径、弦长、圆心角还是运动时间。

\subsection*{模型 1：有界磁场中的圆弧轨迹}

粒子进入有界匀强磁场后，只在磁场内走圆弧。解题关键是找圆心：
\begin{enumerate}
  \item 入射点处，半径垂直于速度方向。
  \item 出射点处，半径也垂直于出射速度方向。
  \item 两条半径的交点就是圆心。
\end{enumerate}

\begin{center}
% TikZ 图：有界磁场圆弧
% Day 22：洛伦兹力
\begin{tikzpicture}[scale=1.22,>=Latex,line cap=round,line join=round]
  \path[use as bounding box] (-1.75,-2.05) rectangle (3.65,2.00);
  \draw[thick,fill=blue!4] (0,-1.48) rectangle (2.75,1.48);
  \foreach \x in {0.45,1.10,1.75,2.40} {
    \foreach \y in {-1.02,-0.34,0.34,1.02} {
      \draw[blue!70!black,thick] (\x-0.055,\y-0.055) -- (\x+0.055,\y+0.055);
      \draw[blue!70!black,thick] (\x-0.055,\y+0.055) -- (\x+0.055,\y-0.055);
    }
  }
  \node[blue!70!black,font=\small] at (2.13,1.13) {$B$ 向里};

  \coordinate (O) at (0,0.35);
  \coordinate (A) at (0,-1.05);
  \coordinate (C) at (1.21,-0.35);
  \draw[->,red,very thick] (-1.30,-1.05) -- (A);
  \node[red,below] at (-0.68,-1.05) {$v_1$};
  \draw[red,very thick,->] (A) arc (-90:-30:1.40);
  \draw[->,red,very thick] (C) -- ++(0.72,1.25);
  \node[red,right] at (1.70,0.58) {$v_2$};

  \fill[red] (A) circle (2pt);
  \node[red,below left=2pt] at (A) {$A$};
  \node[red] at (0.35,-1.30) {$+q$};
  \fill[red] (C) circle (2pt);
  \node[red,below right=2pt] at (C) {$C$};
  \fill (O) circle (1.6pt) node[left=3pt] {$O$};
  \draw[gray,dashed,thick] (O) -- (A);
  \node[gray,left] at (0,-0.35) {$r$};
  \draw[gray,dashed,thick] (O) -- (C);
  \node[gray,above right] at (0.58,0.02) {$r$};
  \node[font=\small] at (1.35,-1.78)
    {入射、出射速度的垂线相交于圆心};
\end{tikzpicture}


\figcaption{有界磁场中先找圆心，再用几何关系求半径和时间}\label{fig:magnetic-boundary-arc}
\end{center}

\textbf{几何细化：} 有界磁场题常把“轨迹圆”藏起来。入射速度和出射速度的垂线交点就是圆心；圆心角 $\alpha$ 等于速度方向转过的角度，时间由圆心角决定：
\[
t=\frac{\alpha}{2\pi}T.
\]
这里最容易错的不是公式，而是把圆心找错。一旦圆心错了，半径、圆心角和运动时间都会跟着错。

\textbf{常见边界关系：} 若粒子从一条边界垂直进入磁场，速度转过 $\alpha$ 后到达与入射边界相距 $d$ 的另一条平行边界，则可由轨迹圆的投影得到
\[
d=r\sin\alpha.
\]
这不是新的物理公式，而是特定边界条件下的几何关系。使用前必须在图中确认 $d$ 对应的是哪一段投影，不能离开图形直接套用。

\subsection*{模型 2：速度选择器}

理想速度选择器中，匀强电场与匀强磁场互相垂直，粒子速度同时垂直于两者。还要选定恰当的传播方向，使电场力与洛伦兹力方向相反。以图中的正电荷向右运动为例，只有满足
\[
|q|E=|q|vB,\qquad \boxed{v=\frac{E}{B}}
\]
的粒子才能沿直线通过。若换成负电荷，两种力同时反向，选出的速率仍为 $E/B$；若把粒子传播方向反过来，磁场力反向而电场力不变，两力便不能平衡。

\begin{center}
% TikZ 图：速度选择器
% Day 22：洛伦兹力
\begin{tikzpicture}[scale=0.92,>=Latex]
  \draw[thick] (-2,1.15) -- (2,1.15);
  \node[above] at (-1.8,1.15) {$+$};
  \draw[thick] (-2,-1.15) -- (2,-1.15);
  \node[below] at (-1.8,-1.15) {$-$};
  \foreach \x in {-1.4,-0.7,0,0.7,1.4} {
    \draw[->,blue,thick] (\x,0.9) -- (\x,-0.9);
  }
  \foreach \x in {-1.5,-0.75,0,0.75,1.5} {
    \foreach \y in {-0.55,0.25} {
      \draw[gray,thick] (\x-0.06,\y-0.06) -- (\x+0.06,\y+0.06);
      \draw[gray,thick] (\x-0.06,\y+0.06) -- (\x+0.06,\y-0.06);
    }
  }
  \fill[red] (-1.8,0) circle (2pt);
  \draw[->,red,very thick] (-1.8,0) -- (1.8,0) node[right] {$v$};
  \draw[->,purple,thick] (-0.3,0) -- (-0.3,-0.85) node[below] {$qE$};
  \draw[->,orange,thick] (0.25,0) -- (0.25,0.85) node[above] {$qvB$};
  \node at (0,-1.65) {$qE=qvB\Rightarrow v=E/B$};
\end{tikzpicture}


\figcaption{速度选择器：图示正电荷沿指定方向运动，两力平衡时不偏转}\label{fig:velocity-selector}
\end{center}

\subsection*{模型 3：质谱仪}

设正离子先通过由匀强电场 $E$ 和匀强磁场 $B_1$ 组成的速度选择器，再进入只有匀强磁场 $B_2$ 的分析区。速度选择器给出
\[
v=\frac{E}{B_1}.
\]
进入分析区后，洛伦兹力提供向心力：
\[
r=\frac{mv}{|q|B_2}.
\]
两式联立可得
\[
\boxed{\frac{m}{|q|}=\frac{B_1B_2r}{E}}.
\]
若粒子走半圆后落到入口边界处的探测器上，入口到落点的距离为 $x=2r$。因此测出落点位置就能确定 $r$，进而得到粒子的质荷比。在 $|q|$、$v$、$B_2$ 相同的情况下，公式中的 $m$ 越大，轨道半径越大。

\begin{center}
\begin{minipage}{0.95\textwidth}
\centering
% TikZ 图：质谱仪原理
% Day 22：洛伦兹力
\begin{tikzpicture}[scale=0.82,>=Latex,line cap=round,line join=round]
  \path[use as bounding box] (-5.10,-1.75) rectangle (3.45,3.55);

  % 速度选择器：E 与 B_1 正交
  \draw[thick] (-4.75,1.15) -- (-1.05,1.15);
  \draw[thick] (-4.75,-1.15) -- (-1.05,-1.15);
  \node at (-4.52,1.42) {$+$};
  \node at (-4.52,-1.42) {$-$};
  \foreach \x in {-4.15,-3.35,-2.55,-1.75} {
    \draw[->,blue,thick] (\x,0.88) -- (\x,-0.88);
  }
  \node[blue,font=\small] at (-1.42,0.84) {$E$ 向下};
  \foreach \x in {-3.82,-3.02,-2.22,-1.42} {
    \foreach \y in {-0.48,0.48} {
      \draw[gray,thick] (\x-0.055,\y-0.055) -- (\x+0.055,\y+0.055);
      \draw[gray,thick] (\x-0.055,\y+0.055) -- (\x+0.055,\y-0.055);
    }
  }
  \node[gray,font=\small] at (-4.18,0.76) {$B_1$ 向里};
  \node[font=\small] at (-2.90,-1.62) {速度选择器：$v=E/B_1$};

  \fill[red] (-4.65,0) circle (2.2pt);
  \node[red,above left] at (-4.65,0) {$+q$};
  \draw[->,red,very thick] (-4.58,0) -- (-0.58,0);
  \node[red,below] at (-2.00,0) {$v$};

  % 分析磁场与探测器
  \draw[thick,fill=blue!4] (-0.55,-0.20) rectangle (3.20,3.20);
  \foreach \x in {0.05,0.75,1.45,2.15,2.85} {
    \foreach \y in {0.35,1.00,1.65,2.30,2.95} {
      \draw[gray] (\x-0.05,\y-0.05) -- (\x+0.05,\y+0.05);
      \draw[gray] (\x-0.05,\y+0.05) -- (\x+0.05,\y-0.05);
    }
  }
  \node[blue!70!black,font=\small] at (2.42,2.90)
    {分析磁场 $B_2$ 向里};

  \coordinate (S) at (-0.55,0);
  \draw[red,very thick,->] (S) arc (-90:90:0.78);
  \draw[orange!90!black,very thick,->] (S) arc (-90:90:1.28);
  \node[red,font=\small] at (0.30,0.80) {$m_1$};
  \node[orange!90!black,font=\small] at (0.86,1.76) {$m_2>m_1$};

  \draw[very thick] (-0.72,0.20) -- (-0.72,2.85);
  \node[font=\small] at (-0.72,3.04) {位置探测器};
  \fill[red] (-0.55,1.56) circle (2pt);
  \fill[orange!90!black] (-0.55,2.56) circle (2pt);
  \draw[->,red,thin] (-1.75,1.56) -- (-0.67,1.56);
  \node[red,left,font=\small] at (-1.75,1.56) {$x_1=2r_1$};
  \draw[->,orange!90!black,thin] (-1.75,2.56) -- (-0.67,2.56);
  \node[orange!90!black,left,font=\small] at (-1.75,2.56) {$x_2=2r_2$};
\end{tikzpicture}


\figcaption{质谱仪原理：速度选择器用 $E,B_1$ 选速，分析磁场 $B_2$ 用落点区分质荷比}\label{fig:mass-spectrometer}
\end{minipage}
\end{center}

\subsection*{模型 4：回旋加速器}

回旋加速器把电场和磁场的作用分开使用：两个 D 形盒内主要有匀强磁场，粒子在盒内做半圆运动；两个盒之间的窄缝中有交变电场，粒子每次经过窄缝时都被加速一次。图中画出正电荷在磁场垂直纸面向里的情形。速度变大后，由
\[
r=\frac{mv}{|q|B}
\]
可知：在 $m$、$|q|$、$B$ 不变时，半径 $r$ 与速度 $v$ 成正比。粒子每过一次缝隙速度变大，下一段圆弧半径也随之变大，所以它不是在同一圆上反复转，而是从中心附近一圈圈向外运动。

\begin{center}
% TikZ 图：回旋加速器模型
% 图示正电荷、磁场垂直纸面向里
\begin{tikzpicture}[scale=0.88,>=Latex,line cap=round,line join=round]
  \path[use as bounding box] (-2.75,-3.25) rectangle (2.75,3.10);

  % 两个 D 形盒与窄缝
  \fill[gray!10] (-0.16,2.20) -- (-0.16,-2.20)
    arc (-90:-270:2.20) -- cycle;
  \fill[gray!10] (0.16,2.20) -- (0.16,-2.20)
    arc (-90:90:2.20) -- cycle;
  \draw[thick] (-0.16,2.20) -- (-0.16,-2.20)
    arc (-90:-270:2.20);
  \draw[thick] (0.16,2.20) -- (0.16,-2.20)
    arc (-90:90:2.20);
  \node[font=\small] at (-1.35,2.52) {D 形盒};
  \node[font=\small] at (1.35,2.52) {D 形盒};

  % 磁场：垂直纸面向里
  \foreach \x in {-1.65,-1.05,-0.52,0.52,1.05,1.65} {
    \foreach \y in {-1.55,-0.78,0.78,1.55} {
      \draw[blue!65!black,thick] (\x-0.06,\y-0.06) -- (\x+0.06,\y+0.06);
      \draw[blue!65!black,thick] (\x-0.06,\y+0.06) -- (\x+0.06,\y-0.06);
    }
  }
  \node[blue!65!black,font=\small] at (0,-2.64)
    {$B$ 垂直纸面向里};

  % 一条连续的、半径逐次增大的轨迹
  \draw[orange!90!black,very thick,->] (0,0) arc (-90:90:0.35);
  \draw[orange!90!black,very thick,->] (0,0.70) arc (90:270:0.65);
  \draw[orange!90!black,very thick,->] (0,-0.60) arc (-90:90:0.95);
  \draw[orange!90!black,very thick,->] (0,1.30) arc (90:270:1.25);
  \draw[orange!90!black,very thick,->] (0,-1.20) arc (-90:66:1.55);
  \fill[red] (0,0) circle (2.3pt);

  % 只画一个电场方向示意，避免把两个时刻画在同一图中
  \draw[->,red,very thick] (-0.48,1.82) -- (0.48,1.82);
  \node[red,right,font=\small] at (0.50,1.82) {缝隙电场 $E$};
  \node[font=\small] at (0,2.86)
    {交变电场每隔 $T/2$ 反向一次};
  \node[font=\small] at (0,-3.02)
    {盒内磁场负责转弯，窄缝电场负责做功};
\end{tikzpicture}


\figcaption{回旋加速器：磁场负责转弯，缝隙中的交变电场负责加速}\label{fig:cyclotron-model}
\end{center}

回旋周期为
\[
T=\frac{2\pi m}{|q|B}
\]
这个式子里没有 $v$，所以在非相对论、匀强磁场的回旋模型中，回旋周期与速度大小无关。这个结论很关键：粒子虽然越来越快，但轨道半径也同步变大，转一圈所需时间仍保持不变。因此只要交变电压频率满足
\[
\boxed{f=\frac{1}{T}=\frac{|q|B}{2\pi m}},
\]
交变电压的周期应等于粒子的回旋周期 $T$，电场方向每经过半个周期 $T/2$ 反向一次。选定合适初相位后，粒子每次到达窄缝时，电场都恰好沿其运动方向，使粒子继续加速。若最大轨道半径为 $R$，则射出速度和最大动能近似为
\[
v_{\max}=\frac{|q|BR}{m},\qquad E_{k\max}=\frac{q^2B^2R^2}{2m}.
\]

\textbf{认知提醒：} 磁场不做功，不能直接增加粒子的动能；真正做功的是窄缝中的电场。磁场的作用是让粒子不断转回窄缝，使电场可以一次次做功。

\section*{六、例题精讲}

\begin{examplebox}{带电粒子的圆周运动}
\textbf{题意分析：} 质子以 $v=2.0\times 10^6\ \mathrm{m/s}$ 的速度垂直进入 $B=0.50\ \mathrm{T}$ 的匀强磁场。质子质量 $m=1.67\times10^{-27}\ \mathrm{kg}$，电荷量 $q=1.60\times10^{-19}\ \mathrm{C}$。求轨道半径。

\textbf{物理模型：} 洛伦兹力提供向心力。

\begin{center}
% TikZ 图：洛伦兹力半径例题
% Day 22：洛伦兹力
\begin{tikzpicture}[scale=0.9,>=Latex]
  \foreach \x in {-1.8,-1.1,-0.4,0.3,1.0,1.7} {
    \foreach \y in {-1.0,-0.35,0.3,0.95} {
      \draw[blue,thick] (\x-0.055,\y-0.055) -- (\x+0.055,\y+0.055);
      \draw[blue,thick] (\x-0.055,\y+0.055) -- (\x+0.055,\y-0.055);
    }
  }
  \node[blue,font=\small] at (2.25,0.95) {$B$ 向里};
  \coordinate (O) at (0,0);
  \coordinate (P) at (0,-1.05);
  \draw[dashed] (O) circle (1.05);
  \draw[red,very thick,->] (-1.45,-1.05) -- (P) node[midway,below] {$v$};
  \draw[red,thick,->] (P) arc (-90:120:1.05);
  \draw[gray] (O) -- (P) node[midway,right] {$r$};
  \fill (O) circle (1.5pt) node[above] {$O$};
  \fill[red] (P) circle (2pt);
\end{tikzpicture}

\end{center}

\textbf{解题过程：}
\[
r=\frac{mv}{|q|B}
=\frac{1.67\times10^{-27}\times2.0\times10^6}
{1.60\times10^{-19}\times0.50}
\approx 4.2\times10^{-2}\ \mathrm{m}
\]

\textbf{方法总结：} 半径公式来自向心力方程，应先理解来源再记忆。速度越大、质量越大，半径越大；电荷量越大、磁场越强，半径越小。
\end{examplebox}

\begin{examplebox}{有界磁场中的圆心、半径与时间}
\textbf{题意分析：} 如图，矩形区域内有垂直纸面向里的匀强磁场，左、右边界相距 $d$。质量为 $m$、电荷量为 $+q$ 的粒子从左边界上的 $A$ 点以速率 $v$ 垂直边界射入，在 $C$ 点离开右边界。出射速度相对入射速度向上偏转 $30^\circ$。求磁感应强度 $B$ 和粒子在磁场中的运动时间。

\textbf{物理模型：} 粒子在磁场内做匀速圆周运动；入射点和出射点处的速度垂线相交于圆心。

\begin{center}
% TikZ 图：有界磁场完整几何例题
% Day 22：洛伦兹力
\begin{tikzpicture}[scale=1.08,>=Latex,line cap=round,line join=round]
  \path[use as bounding box] (-1.35,-0.95) rectangle (4.35,4.35);

  % 矩形有界磁场
  \draw[thick,fill=blue!3] (0,-0.55) rectangle (2.00,4.15);
  \foreach \x in {0.35,1.00,1.65} {
    \foreach \y in {0.15,0.85,1.55,2.25,2.95,3.65} {
      \draw[blue!65!black,thick]
        (\x-0.055,\y-0.055) -- (\x+0.055,\y+0.055)
        (\x-0.055,\y+0.055) -- (\x+0.055,\y-0.055);
    }
  }
  \node[blue!70!black,align=center] at (3.05,2.80)
    {$B$ 垂直纸面\\向里};

  % 几何点：取 r=4，alpha=30 度，横向投影 d=2
  \coordinate (A) at (0,0);
  \coordinate (O) at (0,4);
  \coordinate (C) at (2,0.536);

  % 轨迹、速度和半径
  \draw[red,very thick,->] (A) arc (-90:-60:4);
  \fill[red] (A) circle (2.2pt);
  \fill[red] (C) circle (2.2pt);
  \node[red,left] at (A) {$A$};
  \node[red,right] at (C) {$C$};
  \node[red,below left] at (-0.02,-0.03) {$+q$};
  \draw[->,red,very thick] (A) -- (1.30,0) node[below] {$v$};
  \draw[->,red,very thick] (C) -- ++(1.18,0.68) node[right] {$v_C$};

  \draw[gray!75,dashed,thick] (A) -- (O);
  \draw[gray!75,dashed,thick] (C) -- (O);
  \fill (O) circle (1.8pt);
  \node[above left=1pt] at (O) {$O$};
  \node[left] at (0,2.15) {$r$};
  \node[gray!70!black] at (1.18,2.45) {$r$};
  \draw[orange!90!black,thick,->]
    (0,3.18) arc (-90:-60:0.82);
  \node[orange!90!black,right] at (0.36,3.45) {$30^\circ$};

  % 两边界间距
  \draw[gray,dashed] (0,-0.55) -- (0,-0.85);
  \draw[gray,dashed] (2.00,-0.55) -- (2.00,-0.85);
  \draw[<->,thick] (0,-0.78) -- (2.00,-0.78)
    node[midway,fill=white,inner sep=1pt] {$d$};
  \node[font=\small,rotate=90] at (-0.38,2.10) {左边界};
  \node[font=\small,rotate=90] at (2.34,2.10) {右边界};
\end{tikzpicture}

\end{center}

\textbf{解题过程：}

在 $A$ 点，半径 $OA$ 垂直于入射速度；在 $C$ 点，半径 $OC$ 垂直于出射速度。速度方向转过的角度等于圆心角，因此
\[
\angle AOC=\alpha=30^\circ.
\]
左右边界的距离是半径 $OC$ 在水平方向上的投影：
\[
d=r\sin\alpha=r\sin30^\circ=\frac{r}{2},
\qquad r=2d.
\]
由洛伦兹力提供向心力，
\[
B=\frac{mv}{qr}=\boxed{\frac{mv}{2qd}}.
\]
粒子转过的圆心角为 $\alpha=\pi/6$，故
\[
t=\frac{\alpha}{2\pi}T
=\frac{1}{12}\cdot\frac{2\pi m}{qB}
=\boxed{\frac{\pi m}{6qB}}
=\boxed{\frac{\pi d}{3v}}.
\]

\textbf{方法总结：} 边界题不能从弧线外形直接猜半径。应先在入射点、出射点作速度垂线找圆心，再把边界宽度转化为半径的投影，最后利用圆心角求时间。
\end{examplebox}

\begin{examplebox}{半圆磁场中的运动时间}
\textbf{题意分析：} 一个带正电粒子垂直进入匀强磁场，轨迹刚好为半圆后射出。已知 $m$、$q$、$B$，求粒子在磁场中的运动时间。

\textbf{物理模型：} 圆周运动的一部分。

\begin{center}
% TikZ 图：半圆磁场运动时间
% Day 22：洛伦兹力
\begin{tikzpicture}[scale=0.9,>=Latex]
  \path[use as bounding box] (-2.15,-1.85) rectangle (2.35,1.35);
  \foreach \x in {-1.7,-1.0,-0.3,0.4,1.1,1.8} {
    \foreach \y in {-1.0,-0.35,0.3,0.95} {
      \draw[blue,thick] (\x-0.055,\y-0.055) -- (\x+0.055,\y+0.055);
      \draw[blue,thick] (\x-0.055,\y+0.055) -- (\x+0.055,\y-0.055);
    }
  }
  \node[blue,font=\small] at (2.25,0.95) {$B$ 向里};
  \coordinate (O) at (0,0);
  \coordinate (A) at (0,-1.0);
  \coordinate (C) at (0,1.0);
  \draw[dashed] (O) circle (1.0);
  \draw[red,very thick,->] (-1.35,-1.0) -- (A) node[midway,below] {$v$};
  \draw[red,thick,->] (A) arc (-90:90:1.0);
  \fill[red] (A) circle (2pt) node[below] {入};
  \fill[red] (C) circle (2pt) node[above] {出};
  \fill (O) circle (1.5pt) node[right] {$O$};
  \draw[orange,thick,->] (0,-0.55) arc (-90:90:0.55);
  \node[orange] at (0.72,0) {$\pi$};
\end{tikzpicture}

\end{center}

\textbf{解题过程：}

完整周期为
\[
T=\frac{2\pi m}{|q|B}
\]
半圆对应圆心角 $\pi$，所以时间为半个周期：
\[
t=\frac{T}{2}=\frac{\pi m}{|q|B}
\]

\textbf{方法总结：} 有界磁场时间题看“圆弧对应圆心角”，不是看弧长直观长短。
\end{examplebox}

\begin{examplebox}{速度选择器能通过的速度}
\textbf{题意分析：} 速度选择器中匀强电场强度 $E=3.0\times10^4\ \mathrm{N/C}$，匀强磁场 $B=0.20\ \mathrm{T}$，且 $E\perp B$。粒子沿图示方向进入时，速度同时垂直于 $E$、$B$，电场力和洛伦兹力方向相反。求能沿直线通过的粒子速度。

\textbf{物理模型：} 理想速度选择器中 $E\perp B$，且 $v$ 同时垂直于 $E$、$B$；粒子沿指定方向直线通过意味着横向合力为零。

\begin{center}
% TikZ 图：速度选择器例题
% Day 22：洛伦兹力
\begin{tikzpicture}[scale=0.9,>=Latex,line cap=round]
  \path[use as bounding box] (-2.55,-1.75) rectangle (2.70,1.45);
  \draw[thick] (-2.3,1.0) -- (2.3,1.0);
  \draw[thick] (-2.3,-1.0) -- (2.3,-1.0);
  \node at (-2.15,1.25) {$+$};
  \node at (-2.15,-1.25) {$-$};
  \foreach \x in {-1.5,-0.75,0,0.75,1.5} {
    \draw[->,blue,thick] (\x,0.78) -- (\x,-0.78);
  }
  \foreach \x in {-1.7,-0.9,-0.1,0.7,1.5} {
    \draw[gray,thick] (\x-0.055,-0.055) -- (\x+0.055,0.055);
    \draw[gray,thick] (\x-0.055,0.055) -- (\x+0.055,-0.055);
  }
  \draw[->,red,very thick] (-2.0,0) -- (2.0,0) node[right] {$v$};
  \draw[->,purple,thick] (-0.45,0) -- (-0.45,-0.68);
  \node[purple,font=\small] at (-1.02,-0.46) {$|q|E$};
  \draw[->,orange,thick] (0.45,0) -- (0.45,0.68);
  \node[orange,font=\small] at (0.80,0.55) {$|q|vB$};
  \node[font=\small] at (0,-1.55) {直线通过：$|q|E=|q|vB$};
\end{tikzpicture}

\end{center}

\textbf{解题过程：}
\[
|q|E=|q|vB.
\]
两边都有 $|q|$，所以在这个理想正交场模型中，能通过的速率与粒子电荷量大小和正负均无关：
\[
v=\frac{E}{B}
=\frac{3.0\times10^4}{0.20}
=1.5\times10^5\ \mathrm{m/s}.
\]

\textbf{方法总结：} 速度选择器题不能只背 $v=E/B$，还要先检查 $E\perp B$、$v\perp E$、$v\perp B$，并确认两种力方向相反。条件满足后，速度过大或过小都会偏转。
\end{examplebox}

\begin{examplebox}{质谱仪中的半径比较}
\textbf{题意分析：} 两种一价正离子经过同一速度选择器后进入同一匀强磁场。若质量分别为 $m$ 和 $2m$，比较它们在磁场中的轨道半径。

\textbf{物理模型：} 速度相同、电荷量相同、磁场相同，半径只随质量变化。

\begin{center}
% TikZ 图：质谱仪半径比较例题
% Day 22：洛伦兹力
\begin{tikzpicture}[scale=0.9,>=Latex]
  \draw[thick,fill=blue!4] (-1.4,-1.55) rectangle (2.6,1.55);
  \foreach \x in {-0.9,-0.2,0.5,1.2,1.9} {
    \foreach \y in {-0.95,-0.25,0.45,1.15} {
      \draw[gray] (\x-0.05,\y-0.05) -- (\x+0.05,\y+0.05);
      \draw[gray] (\x-0.05,\y+0.05) -- (\x+0.05,\y-0.05);
    }
  }
  \node[blue,font=\small] at (2.1,1.25) {$B$ 向里};
  \draw[->,red,very thick] (-2.15,0) -- (-1.4,0) node[midway,below] {$v$};
  \draw[red,thick,->] (-1.4,0) arc (180:15:0.72);
  \draw[orange,thick,->] (-1.4,0) arc (180:15:1.22);
  \node[red] at (-0.65,0.92) {$m$};
  \node[orange] at (0.10,1.48) {$2m$};
  \node[font=\small] at (0.55,-1.85) {$r_2=2r_1$};
\end{tikzpicture}

\end{center}

\textbf{解题过程：}
\[
r=\frac{mv}{|q|B}.
\]
对质量为 $m$ 的离子：
\[
r_1=\frac{mv}{|q|B}.
\]
对质量为 $2m$ 的离子：
\[
r_2=\frac{2mv}{|q|B}=2r_1.
\]

\textbf{方法总结：} 质谱仪题先确认粒子速度和电荷量是否相同，再用 $r=\dfrac{mv}{|q|B}$ 比较半径；在本题条件下，半径与质量成正比。
\end{examplebox}

\section*{七、变式训练}

\begin{enumerate}
  \item 正电荷和负电荷以相同速度水平向右进入垂直纸面向里的匀强磁场。分别判断二者刚进入磁场时所受洛伦兹力方向，并说明轨迹向哪一侧弯曲。
  \item 电子以速度 $v$ 垂直进入匀强磁场，若磁场增强为原来的 2 倍，轨道半径如何变化？
  \begin{center}
  \begin{tikzpicture}[scale=0.88,>=Stealth,line cap=round,line join=round]
  \path[use as bounding box] (-0.35,-0.70) rectangle (7.25,1.98);
  \begin{scope}
    \fill[blue!5] (-0.15,-0.1) rectangle (2.55,1.95);
    \foreach \x in {0.25,0.75,1.25,1.75,2.25}
      \foreach \y in {0.25,0.75,1.25,1.75}
        \node[blue!70] at (\x,\y) {$\times$};
    \coordinate (S) at (0.42,1.72);
    \coordinate (O) at (0.42,0.62);
    \draw[dashed] (O) -- (S) node[midway,left,font=\small] {$r$};
    \fill (O) circle (0.025);
    \draw[thick,red,->] (S) arc[start angle=90,end angle=-35,radius=1.10];
    \draw[->,thick] (-0.22,1.72) -- (S) node[midway,above,font=\small] {$v$};
    \node[red,font=\small] at (1.12,1.15) {$e^-$};
    \node[font=\small] at (1.20,-0.42) {$B$ 向里};
  \end{scope}
  \draw[->,thick] (2.78,0.88) -- (3.28,0.88);
  \node[font=\scriptsize] at (3.03,1.17) {$B$ 加倍};
  \begin{scope}[xshift=3.55cm]
    \fill[blue!5] (-0.15,-0.1) rectangle (2.55,1.95);
    \foreach \x in {0.25,0.75,1.25,1.75,2.25}
      \foreach \y in {0.25,0.75,1.25,1.75}
        \node[blue!80] at (\x,\y) {$\times$};
    \coordinate (S2) at (0.42,1.72);
    \coordinate (O2) at (0.42,1.17);
    \draw[dashed] (O2) -- (S2) node[midway,left,font=\small] {$r/2$};
    \fill (O2) circle (0.025);
    \draw[thick,red,->] (S2) arc[start angle=90,end angle=-70,radius=0.55];
    \draw[->,thick] (-0.22,1.72) -- (S2) node[midway,above,font=\small] {$v$};
    \node[red,font=\small] at (0.95,0.95) {$e^-$};
    \node[font=\small] at (1.20,-0.42) {$2B$ 向里};
  \end{scope}
\end{tikzpicture}

  \end{center}
  \item 一个带电粒子在匀强磁场中运动半径为 $0.20\ \mathrm{m}$，速度为 $1.0\times10^6\ \mathrm{m/s}$。若 $B=0.10\ \mathrm{T}$，求其比荷大小 $|q|/m$。
  \item 粒子在磁场中偏转 $90^\circ$ 后射出，求其在磁场中运动时间与周期的关系。
  \begin{center}
  % TikZ 图：变式训练中的四分之一圆轨迹
% Day 22：洛伦兹力
\begin{tikzpicture}[scale=1.02,>=Latex,line cap=round,line join=round]
  \path[use as bounding box] (-1.65,-1.55) rectangle (2.80,1.45);
  \draw[thick,fill=blue!4] (0,-1.35) rectangle (1.65,1.25);
  \foreach \x in {0.38,0.92,1.46} {
    \foreach \y in {-0.85,-0.25,0.35,0.95} {
      \draw[blue!70!black,thick] (\x-0.05,\y-0.05) -- (\x+0.05,\y+0.05);
      \draw[blue!70!black,thick] (\x-0.05,\y+0.05) -- (\x+0.05,\y-0.05);
    }
  }
  \node[blue!70!black,font=\small] at (2.30,0.62) {$B$ 向里};

  \coordinate (O) at (0,0);
  \coordinate (A) at (0,-1);
  \coordinate (C) at (1,0);
  \draw[->,red,very thick] (-1.25,-1) -- (A);
  \node[red] at (-0.72,-1.20) {$v_1$};
  \draw[red,very thick,->] (A) arc (-90:0:1);
  \draw[->,red,very thick] (C) -- ++(0,0.88);
  \node[red,right] at (1.10,0.78) {$v_2$};

  \draw[gray,dashed,thick] (O) -- (A);
  \draw[gray,dashed,thick] (O) -- (C);
  \fill (O) circle (1.5pt) node[above left] {$O$};
  \fill[red] (A) circle (2pt);
  \node[red] at (-0.12,-1.30) {$A$};
  \node[red] at (0.34,-1.22) {$+q$};
  \fill[red] (C) circle (2pt) node[below right=2pt] {$C$};
  \draw[orange,thick,->] (0,-0.45) arc (-90:0:0.45);
  \node[orange] at (0.35,-0.34) {$90^\circ$};
\end{tikzpicture}

  \end{center}
  \item 正离子从左边界垂直进入匀强磁场，速度方向转过 $30^\circ$ 时恰好到达与左边界相距 $d$ 的右边界。求轨道半径，并写出它在磁场中的运动时间与周期的关系。
  \item 速度选择器中 $E=2.0\times10^4\ \mathrm{N/C}$，$B=0.20\ \mathrm{T}$，能直线通过的粒子速度是多少？
  \item 两种一价正离子经同一速度选择器后进入同一分析磁场，测得半圆轨迹直径之比为 $x_1:x_2=3:5$。求两种离子的质量之比。
  \item 某回旋加速器的磁感应强度为 $B$，D 形盒最大半径为 $R$。粒子的质量为 $m$、电荷量大小为 $|q|$。写出交变电压的频率及粒子射出时的最大动能。
\end{enumerate}

\begin{learningbox}{变式训练参考结果}
\begin{enumerate}[itemsep=0.45em,topsep=0.3em,parsep=0pt]
  \item 对正电荷，$\vect{v}\times\vect{B}$ 方向向上，轨迹向上弯曲；负电荷所受力方向相反，轨迹向下弯曲。
  \item 由 $r=mv/(|q|B)$，磁感应强度变为原来的 2 倍时，轨道半径变为原来的一半。
  \item $\displaystyle \frac{|q|}{m}=\frac{v}{Br}
  =5.0\times10^7\ \mathrm{C/kg}$。
  \item 偏转 $90^\circ$ 对应转过四分之一圆周，故在磁场中的运动时间为 $T/4$。
  \item 由 $d=r\sin30^\circ$ 得 $r=2d$；偏转 $30^\circ$ 对应六分之一圆周，故 $t=T/12$。
  \item 直线通过时电场力与洛伦兹力平衡，故
  $v=E/B=1.0\times10^5\ \mathrm{m/s}$。
  \item 同一速度选择器和分析磁场中，$x=2r\propto m/|q|$。两种离子的 $|q|$ 相同，故 $m_1:m_2=3:5$。
  \item $\displaystyle f=\frac{|q|B}{2\pi m}$，
  $\displaystyle E_{k\max}=\frac{q^2B^2R^2}{2m}$。
\end{enumerate}
\end{learningbox}

\section*{八、竞赛入门补充}

\begin{competitionbox}
\textbf{螺旋线运动：速度不完全垂直磁场}

若粒子速度与磁场有夹角，可以把速度分解为
\[
v_\parallel = v\cos\theta,\qquad v_\perp = v\sin\theta
\]

\begin{itemize}
  \item 平行磁场的分量 $v_\parallel$ 不受洛伦兹力影响，保持匀速直线运动。
  \item 垂直磁场的分量 $v_\perp$ 产生圆周运动。
  \item 两个运动合成后，轨迹为螺旋线。
\end{itemize}

\begin{center}
% TikZ 图：斜入磁场的立体螺旋运动
% Day 22：洛伦兹力
\begin{tikzpicture}[
  scale=0.9,
  >=Latex,
  line cap=round,
  line join=round,
  x={(1cm,0cm)},
  y={(0cm,1cm)},
  z={(0.34cm,0.20cm)}
]
  \path[use as bounding box] (-3.65,-1.85) rectangle (3.85,1.80);

  % 左侧：入射速度分解
  \begin{scope}[shift={(-3.15,0)}]
    \coordinate (P) at (0,0);
    \fill[red] (P) circle (2pt);
    \draw[->,orange!90!black,very thick] (P) -- (1.05,0.82);
    \node[orange!90!black,above] at (0.72,0.60) {$v$};
    \draw[->,orange!90!black,thick] (P) -- (1.05,0);
    \node[orange!90!black,below] at (0.55,0) {$v_\parallel$};
    \draw[->,orange!90!black,thick] (1.05,0) -- (1.05,0.82);
    \node[orange!90!black,right] at (1.05,0.42) {$v_\perp$};
  \end{scope}

  % 右侧：沿磁场方向推进的立体螺旋线
  \draw[->,blue!70!black,very thick]
    (-1.45,0,0) -- (3.15,0,0)
    node[right] {$B$};
  \node[blue!70!black,above] at (1.05,0.72,0)
    {螺旋轴线 $\parallel B$};

  % 起、末截面帮助读者看出轨迹绕轴旋转
  \draw[gray!55,thin,domain=0:360,samples=80,variable=\t]
    plot ({-1.25},{0.52*cos(\t)},{0.52*sin(\t)});
  \draw[gray!55,thin,domain=0:360,samples=80,variable=\t]
    plot ({2.49},{0.52*cos(\t)},{0.52*sin(\t)});

  \draw[red,very thick,->,domain=0:720,samples=220,variable=\t]
    plot ({-1.25+0.0052*\t},{0.52*cos(\t)},{0.52*sin(\t)});

  % 半径与螺距
  \draw[->,orange!90!black,thick]
    (0.62,0,0) -- (0.62,0.52,0)
    node[midway,left] {$r$};
  \draw[<->,purple,thick]
    (-0.32,-0.85,0) -- (1.55,-0.85,0)
    node[midway,below] {$h=v_\parallel T$};
  \draw[gray,dashed] (-0.32,-0.72,0) -- (-0.32,0.52,0);
  \draw[gray,dashed] (1.55,-0.72,0) -- (1.55,0.52,0);

  \node[font=\small] at (0,-1.55)
    {沿磁场方向匀速前进 $+$ 垂直磁场方向圆周运动 $=$ 螺旋运动};
\end{tikzpicture}


\figcaption{速度斜入磁场：垂直分量形成圆周运动，平行分量使轨迹沿磁场方向推进}\label{fig:helical-motion}
\end{center}

螺旋线半径和螺距分别为
\[
r=\frac{mv_\perp}{|q|B},\qquad h=v_\parallel T
\]
这个模型是理解带电粒子在地磁场中运动、等离子体约束和粒子探测器轨迹的入门基础。

\textbf{严谨提醒：} 螺旋运动中磁场仍然不做功，总速率不变；只是把运动分解后，平行分量保持匀速，垂直分量不断改变方向。
\end{competitionbox}

\subsection*{1.\quad 霍尔效应：用洛伦兹力测磁场}

霍尔效应是洛伦兹力在现代传感器中的典型应用。导体或半导体中有电流通过时，载流子具有定向速度；若再加一个垂直于电流方向的磁场，载流子受到洛伦兹力向一侧偏转，导体两侧积累异号电荷，形成横向电势差，叫\textbf{霍尔电压}。下图按正载流子模型作图：正载流子向上偏转，上侧积累正电荷。若实际载流子是电子，在约定电流和磁场方向不变时，电子漂移方向与约定电流相反，负电荷又使洛伦兹力方向反向；电子仍向图中的上侧偏转，但上侧积累负电荷，因此霍尔电压的极性与正载流子情形相反。

\begin{center}
% TikZ 图：霍尔效应的正载流子等效模型
\begin{tikzpicture}[scale=1.05,>=Stealth,line cap=round]
  \path[use as bounding box] (-3.35,-1.95) rectangle (3.55,1.95);

  \draw[rounded corners=3pt,thick,fill=blue!4]
    (-2.40,-0.88) rectangle (2.40,0.88);
  \foreach \x in {-1.75,-0.95,-0.15,0.65,1.45} {
    \node[font=\small,blue!70!black] at (\x,0.55) {$+$};
    \node[font=\small,blue!70!black] at (\x,-0.55) {$-$};
  }

  \draw[->,very thick,red] (-2.95,0) -- (2.95,0);
  \node[red,above,font=\small] at (1.88,0) {$v_d,\ I$};
  \node[red,font=\small] at (0,-1.30)
    {正载流子等效：$v_d$ 与约定电流 $I$ 同向};

  \draw[->,orange!85!black,very thick] (-1.32,0) -- (-1.32,0.75);
  \node[orange!85!black,left,font=\small] at (-1.32,0.38) {$F_L$};
  \draw[->,purple,very thick] (0.30,0.52) -- (0.30,-0.52);
  \node[purple,right,font=\small] at (0.30,0) {$E_H$};

  \draw[<->,thick,gray] (-2.72,-0.84) -- (-2.72,0.84);
  \node[gray,left] at (-2.72,0) {$d$};
  \draw[<->,thick,blue!70!black] (2.78,-0.80) -- (2.78,0.80);
  \node[blue!70!black,right,font=\small] at (2.78,0) {$U_H$};
  \node[blue!70!black] at (3.18,0.64) {$+$};
  \node[blue!70!black] at (3.18,-0.64) {$-$};

  \foreach \x in {-1.80,-0.90,0,0.90,1.80} {
    \node[font=\large] at (\x,1.34) {$\times$};
    \node[font=\large] at (\x,-1.76) {$\times$};
  }
  \node[font=\small] at (0,1.72) {$B$ 垂直纸面向里};
\end{tikzpicture}


\figcaption{霍尔效应：载流子受洛伦兹力偏转，横向电荷分离形成霍尔电压}\label{fig:hall-effect-model}
\end{center}

随着载流子在横向积累，横向电场逐渐增强。达到稳定状态时，横向电场力与洛伦兹力大小相等：
\[
|q|E_H=|q|v_dB.
\]
若导体横向宽度为 $d$，两侧霍尔电压大小为 $|U_H|$，则在匀强横向电场近似下
\[
E_H=\frac{|U_H|}{d},\qquad
\boxed{|U_H|=v_dBd}.
\]
设导体沿磁场方向的厚度为 $h$，单位体积载流子数为 $n$，则电流
\[
I=n|q|v_d(dh).
\]
消去 $v_d$ 可得霍尔电压大小
\[
\boxed{|U_H|=\frac{BI}{n|q|h}}.
\]
测量霍尔电压的大小可以反求磁感应强度，测量其极性还可以判断主要载流子的正负。霍尔传感器据此用于磁场测量、无接触测速和电机转角检测。解题时要先确定载流子速度方向，再由电荷正负判断偏转方向，最后写稳定时的横向平衡条件。

\subsection*{2.\quad 补充训练：磁场边界题的三问}

遇到有界磁场，先连续回答三问：\textbf{入射点和出射点在哪里？这两点处速度方向分别是什么？由速度垂线找到的圆心在哪里？} 圆心确定后，半径、圆心角和时间才有可靠的几何依据；不要在圆心尚未确定时直接套用半径或周期公式。

\section*{九、易错点}

\begin{mistakebox}
\begin{enumerate}
  \item \textbf{负电荷方向忘记反向：} 先按正电荷用左手定则判断，再把方向反过来给负电荷。
  \item \textbf{把“有加速度”和“速率增大”混为一谈：} 判断速率是否变化时，先看力在速度方向上有没有分量。纯磁场中洛伦兹力始终垂直速度，故只改变方向、提供向心加速度，不单独改变速率。
  \item \textbf{把安培力和洛伦兹力割裂开：} 在高中采用的正载流子等效模型中，安培力可以理解为载流子所受洛伦兹力经导体内部相互作用传递后形成的宏观表现。
  \item \textbf{把半径公式写成 $r=\dfrac{qB}{mv}$：} 先由 $|q|vB=m\dfrac{v^2}{r}$ 推出，再写 $r=\dfrac{mv}{|q|B}$。
  \item \textbf{周期误以为与速度有关：} 先由 $T=2\pi r/v$ 与 $r=mv/(|q|B)$ 联立，观察 $v$ 被约去，才得到 $T=\dfrac{2\pi m}{|q|B}$。这一结论限于非相对论、匀强磁场且速度垂直磁场的圆周运动模型。
  \item \textbf{边界题不找圆心：} 先在入射点、出射点作速度垂线找圆心，再由半径和圆心角列几何关系。
\end{enumerate}
\end{mistakebox}

\section*{十、今日小结}

\begin{enumerate}
  \item 洛伦兹力大小 $F=|q|vB\sin\theta$，列式前先判断 $v$ 与 $B$ 的夹角；判方向时还要检查电荷正负。
  \item 在正载流子等效模型中，安培力可理解为洛伦兹力经导体传递后的宏观表现；磁场洛伦兹力始终垂直速度，因此不做功、不改变速率。
  \item 垂直进入匀强磁场时，洛伦兹力提供向心力，轨道半径 $r=\dfrac{mv}{|q|B}$，周期 $T=\dfrac{2\pi m}{|q|B}$。
  \item 有界磁场题先由入射、出射速度的垂线找圆心，再用边界几何求半径、圆心角和时间。
  \item 速度选择器负责选速，质谱仪利用轨迹测质荷比，回旋加速器用电场加速、磁场转弯；霍尔效应则把载流子的横向偏转转化为可测电压。
\end{enumerate}

\textbf{一分钟复盘：} 画一个带电粒子垂直进入匀强磁场的图，标出 $v$、$B$、$F$、圆心和半径，并分别判断正、负电荷的偏转方向；最后检查半径是否垂直速度、洛伦兹力是否只改变方向而不改变速率。
